\documentclass[11pt]{article}
\usepackage[a4paper,margin=1in]{geometry}
\usepackage{amsmath,amssymb,amsthm,booktabs,graphicx,hyperref,natbib}
\hypersetup{hidelinks}
\newtheorem{theorem}{Theorem}
\newtheorem{proposition}{Proposition}
\newtheorem{corollary}{Corollary}
\newtheorem{definition}{Definition}
\title{Auditing Ecological State Classifications after Structural Change:\\Decision Sufficiency, Minimal Repair, and Historical Context}
\author{Author names withheld for drafting}
\date{}

\begin{document}
\maketitle

\begin{abstract}
Ecological management often reuses coarse state classifications after species turnover, interaction rewiring, habitat change, or changes in the available interventions. A classification that was adequate when it was defined can then become insufficient for a new management decision even if all model parameters are known perfectly. We formalize this problem as an audit of the decision sufficiency of an inherited ecological state variable. A source classification is carried to a target system through a declared structural relation and tested against target outputs, legal actions, and action-conditioned successors. Failure has a local operational witness. We then apply established coarsest-partition refinement machinery to return the least exact distinction that must be added to the inherited interface; the algorithm itself is not new. We separate the resulting structural repair complexity from distribution-sensitive information, monitoring cost, and decision regret, and show how separating witnesses define requirements for targeted monitoring. For multiple declared replacement routes, we characterize when one route-independent carried interface exists and when immutable historical context is necessary. A finite plant--pollinator example, supplemented by an illustrative probabilistic decision analysis, shows how an inherited functional class can reverse restoration priority after turnover. The framework is therefore a structural precondition for adaptive management: before optimizing a policy or collecting more data, it asks whether the state representation supplied to that decision still contains the distinctions that the target intervention can expose.
\end{abstract}

\noindent\textbf{Keywords:} ecological state classification; adaptive management; structural change; targeted monitoring; model transfer; decision sufficiency

\section{Introduction}
Ecological management depends on compressed descriptions of complex systems. Vegetation states, habitat classes, functional guilds, occupancy states, resilience categories, and other macrostates allow monitoring and intervention to proceed without representing every species, interaction, demographic coordinate, and historical detail. State-and-transition models make this compression explicit by linking ecological states to expected pathways of change and management actions, and empirical work has shown how such models can support targeted monitoring by identifying a reduced set of variables that distinguish predefined states and transitions \citep{JonesEtAl2023}. Adaptive-management frameworks likewise require a state representation on which action consequences can be evaluated \citep{Holling1978,Walters1986}.

The difficulty is that the adequacy of a state classification need not survive ecological change. Species can disappear or colonize, interaction networks can rewire, habitats can be replaced, and management can introduce interventions that were unavailable when the classification was designed. Two configurations that were safely merged under the source regime can then differ in an output that matters, in which actions are feasible, or in where the same action leads. In that case the problem is not merely that a transition probability or regression coefficient needs to be re-estimated. The inherited state variable itself suppresses a distinction required by the target decision.

Ecologists have identified transferability as a persistent challenge for prediction across places and conditions \citep{YatesEtAl2018}. Most familiar transfer problems are statistical: whether fitted relationships remain accurate, how extrapolation should be diagnosed, or how uncertainty should be propagated. The failure considered here is deliberately prior to those questions. It can occur with perfect knowledge of both source and target systems and with no sampling or parameter error. The question is structural:
\begin{quote}
When does an ecological state classification that was accepted for a source management problem remain sufficient after structural change, and what is the least additional distinction required when it does not?
\end{quote}

Formal work already supplies much of the mathematical machinery needed to answer parts of this question. Lumpability and bisimulation characterize exact aggregation within a specified process \citep{KemenySnell1960,FeretEtAl2012,LarsenSkou1991}. Coarsest stable refinement from an arbitrary initial partition is established in partition-refinement theory and MDP model minimization \citep{PaigeTarjan1987,GivanDeanGreig2003}. State abstraction has been systematized by what it preserves \citep{LiWalshLittman2006}, maps between distinct decision processes are treated by homomorphism theory \citep{RavindranBarto2003}, and state abstractions have been explicitly transferred from source MDPs to target MDPs \citep{WalshLiLittman2006Transfer}. We therefore claim none of these constructions as new.

The ecological gap is narrower. Managers often inherit a classification that already has an accepted interpretation and a monitoring infrastructure. After a declared structural change, the relevant question is not simply how to compute another efficient abstraction of the target. It is whether that inherited interface still supports the target actions, where it fails if it does not, and which distinctions a monitoring design must recover before the old decision rule can be trusted. This differs from targeted monitoring that begins with predefined states and asks which measured variables best discriminate them \citep{JonesEtAl2023}: here the first question is whether the predefined states themselves remain decision-sufficient. It also differs from value-of-information analysis, which asks whether reducing uncertainty is worth the cost given objectives, actions, and consequences \citep{CanessaEtAl2015}; our audit can identify a structurally missing distinction before uncertainty about that distinction is assigned a probability or a sampling design.

We organize these ingredients into a source-relative ecological audit. A source macro-law is carried to a target system through a declared relation between state spaces that need not coincide. The inherited target fibers are tested for equality of current output, legal-action rows, and successor labels. If they pass, the old state variable remains exact for the declared target interface. If they fail, a pair of states provides a local obstruction, and established coarsest-refinement machinery identifies the least exact split of the inherited classes. That split can be read as a monitoring requirement: candidate measurements must separate the state pairs that the old classification merges but the repaired interface distinguishes.

We add two further pieces needed for ecological interpretation. First, we separate four consequences that should not be conflated: structural repair complexity, distribution-sensitive additional information, the cost of realizing the distinction with measurements, and decision regret. State-count defect alone is not an ecological effect size and does not determine management loss. Second, for systems connected by several declared replacement routes, we ask whether alternative histories induce the same terminal carried interface. History-sensitive behavioral equivalences are well established in formal models of concurrency \citep{MontanariPistore1997}; our narrower object is the externally declared replacement route and the source label map that it carries. We characterize when route identity can be forgotten and when at least one immutable context per distinct carried terminal map is necessary.

The result is best viewed as a theory of ecological state representation for management under structural change. It does not optimize management policies, infer source--target correspondences, or estimate stochastic dynamics. Instead it audits a prerequisite for those tasks: whether the ecological state variable supplied to prediction and intervention still contains the distinctions that the target management problem can operationally expose.

\section{Relation to existing theory and practice}
Table~\ref{tab:boundaries} summarizes the main boundaries. The novelty claim is intentionally modest: established abstraction and refinement results are used directly; the contribution is their organization around an inherited ecological decision interface, the translation of failure into a monitoring requirement, and the route-dependent completion problem.

\begin{table}[t]
\centering
\small
\begin{tabular}{p{0.23\linewidth}p{0.31\linewidth}p{0.38\linewidth}}
\toprule
Literature & Established question & Role or boundary here \\
\midrule
Lumpability and bisimulation & Which states can be aggregated while preserving specified behavior? & Supplies exactness concepts for a declared system. \\
Partition refinement / MDP minimization & What is the coarsest stable refinement of a given partition? & Supplies the repair construction; not claimed as new. \\
State-abstraction transfer and homomorphisms & How can abstractions or dynamics-preserving maps be reused across decision processes? & Shows that source-to-target abstraction transfer itself is prior art; we audit an inherited ecological interface and treat failure as the object of interest. \\
Ecological transferability & When do ecological predictions generalize across settings? & Establishes the importance of transfer; our failure is structural rather than statistical. \\
State-and-transition monitoring & Which variables and thresholds distinguish predefined ecological states and transitions? & We ask first whether the predefined states remain sufficient after the action repertoire or structure changes. \\
Value of information & Is additional information worth collecting for a decision under uncertainty? & Operates downstream of our structural audit once candidate missing distinctions and decision consequences are defined. \\
History-preserving formalisms & When must causal execution history be retained in behavioral equivalence? & Related history-sensitive machinery; our scoped result concerns externally declared replacement routes and carried terminal label maps. \\
\textbf{Present framework} & \textbf{Can an inherited ecological state interface still support the target management problem?} & \textbf{Return pass/fail, a local witness, the least exact distinction, a monitoring target, and required route context.} \\
\bottomrule
\end{tabular}
\caption{Boundary between neighboring theories and the ecological decision-sufficiency audit developed here}
\label{tab:boundaries}
\end{table}

\subsection{Exact aggregation and refinement}
Lumpability and behavioral equivalence provide exact standards for when a coarse variable factors the behavior of a fixed process \citep{KemenySnell1960,FeretEtAl2012,LarsenSkou1991}. The construction used below for repair is likewise established. The relational coarsest-partition problem accepts an arbitrary starting partition \citep{PaigeTarjan1987}, and MDP model minimization computes the coarsest homogeneous refinement of any partition in a stochastic setting that contains our deterministic case \citep{GivanDeanGreig2003}. Consequently, finite termination, coarseness, and uniqueness of that refinement are not Paper A discoveries.

\subsection{Abstraction transfer between systems}
Source-to-target reuse is also not new in itself. Homomorphism theory relates distinct decision processes under dynamics-preserving maps \citep{RavindranBarto2003}, and \citet{WalshLiLittman2006Transfer} explicitly develop an algorithm for transferring state abstractions learned from source MDPs to target MDPs. Their objective is efficient planning and learning in target environments under transferable abstraction structure. Our ecological question is different but narrower: a classification with an accepted source meaning is carried through a declared ecological correspondence, and we ask whether the target management repertoire falsifies that inherited merge. The correspondence need not be a valid homomorphism; detecting its failure with respect to the inherited interface is precisely the audit.

\subsection{Ecological state models, monitoring, and decision analysis}
State-and-transition models provide a natural applied setting because they connect ecological state labels to transitions, drivers, and management expectations. \citet{JonesEtAl2023} show how field data can validate expert-defined vegetation states and identify a small number of variables and thresholds for targeted monitoring. We address the complementary upstream problem: structural change may make two configurations within one inherited state respond differently to a target intervention, so the state definition itself must first be audited. Once a separating ecological distinction has been identified, methods such as those of \citet{JonesEtAl2023} can be used to find practical indicators and thresholds.

Decision analysis addresses a different downstream question. Value-of-information methods quantify whether collecting additional data is worth its cost given uncertainty, objectives, actions, and consequences \citep{CanessaEtAl2015}. The present framework does not calculate VoI. It identifies when an inherited interface suppresses a distinction that makes action consequences nonuniform. That distinction can then become a target for measurement design and, where uncertainty matters, for a VoI calculation.

\subsection{Historical context}
Ecological legacy effects motivate retaining disturbance, colonization, or management history, but retaining all history can make a state description unnecessarily large. Formal computer-science literatures likewise contain history-preserving equivalences that retain causal information about executions \citep{MontanariPistore1997}. We do not claim invention of history-sensitive state descriptions. Our scoped question is whether several declared structural-replacement routes carry the same source label assignment to a shared terminal system. If they do, route identity can be discarded for this inherited interface; if they do not, a route-free carried label does not exist.

\section{Operational framework}
A finite controlled ecological system is a tuple
\[
\mathcal{S}=(X,A,G,\delta,y),
\]
where $X$ is the state set, $A$ is the declared action set, $G$ is a finite legal-action grammar, $\delta$ is the controlled transition rule, and $y$ is the observable output. A macro-law is a surjection $\pi:X\times G\to Q$.

\begin{definition}[Exact operational macro-law]
A macro-law is exact when states in the same macro-fiber have identical current outputs, identical legal-action rows, and successors carrying the same macro-label under every declared legal action.
\end{definition}

Structural change from a source system $\mathcal{S}_0$ to a target system $\mathcal{S}_1$ is represented by a declared total relation
\[
R\subseteq (X_0\times G_0)\times (X_1\times G_1).
\]
When the relation is source-label consistent and covers the target, the source law induces a carried target classification $C=\pi_R$. The relation is treated as declared scientific input; this paper does not infer it from data.

\section{Results}
\subsection{Decision-sufficiency audit}
\begin{theorem}[Operational portability]
For a finite exact source macro-law and a declared label-consistent structural-change relation, the carried target classification is exact if and only if current outputs, legal-action rows, and successor carried labels are constant inside every carried target fiber.
\end{theorem}

This criterion is the audit: it asks whether the old state variable still contains every distinction required by the declared target outputs and interventions. Failure is structural and therefore cannot be repaired merely by re-estimating parameters while keeping the same state partition.

\begin{proposition}[Local obstruction]
Failure of portability is witnessed by two target states in one inherited fiber that differ in current output, legal actions, or successor inherited labels under a declared action or finite legal future.
\end{proposition}

A witness is useful ecologically because it identifies a concrete contrast that the inherited state variable suppresses. If the difference is caused, for example, by substitute-pollinator access, soil condition, or disturbance legacy, candidate monitoring variables must be capable of separating the witness states before the inherited management rule can be restored.

\subsection{Least exact distinction by established refinement}
Let $Q^\ast$ denote the coarsest exact target partition refining the carried classification $C$.

\begin{proposition}[Source-relative exact repair]
Applying the established coarsest stable refinement construction to $C$ yields $Q^\ast$, the unique coarsest exact target partition that never merges states separated by the inherited classification.
\end{proposition}

This proposition is an instantiation of established results \citep{PaigeTarjan1987,GivanDeanGreig2003}, not a new generic refinement theorem. Its ecological role is to define a canonical answer to a specific management question: which distinctions must be added while preserving every inherited merge that remains compatible with the target interface?

\subsection{Separate structural, informational, monitoring, and decision consequences}
The repair has several consequences, but they should not be collapsed into one number.

\begin{definition}[Structural transport defect]
Let $C$ be the carried target classification and $Q^\ast$ its coarsest exact refinement. Define
\[
\Delta_{\#}=|Q^\ast|-|C|,
\qquad
\Delta_{K}=\log_2|Q^\ast|-\log_2|C|.
\]
\end{definition}

$\Delta_{\#}$ records the increase in the number of operational states and $\Delta_K$ is a uniform code-length proxy. Neither quantity is, by itself, a monitoring cost, an ecological effect size, or a measure of management regret. In the verified accumulating family with $m$ independently exposed binary distinctions,
\[
|C|=2,\qquad |Q^\ast|=2^m+1,
\]
so $\Delta_{\#}=2^m-1$. This is a sharp finite witness about loss of compression, not a universal ecological scaling law.

If a target-state distribution $\mu$ is available, the additional label information can be represented separately.

\begin{definition}[Conditional repair information]
Treat $C$ and $Q^\ast$ as the carried and repaired labels of a target microstate drawn from $\mu$. Define
\[
D_H(\mu)=H_{\mu}(Q^\ast\mid C).
\]
\end{definition}

\begin{proposition}[Minimum conditional label information]
For any exact target refinement $P$ of $C$,
\[
H_{\mu}(Q^\ast\mid C)\le H_{\mu}(P\mid C).
\]
If $k_c$ repaired blocks occur inside carried class $c$, then
\[
0\le D_H(\mu)\le \sum_c \mu(c)\log_2 k_c\le \log_2\max_c k_c.
\]
\end{proposition}

The result follows because every exact refinement of $C$ refines the coarsest exact repair $Q^\ast$. Unlike $\Delta_K$, $D_H$ depends on how frequently the target states are expected to occur.

The abstract repair also induces a concrete monitoring-design problem. Let
\[
E=\{(x,x'):C(x)=C(x'),\;Q^\ast(x)\ne Q^\ast(x')\}
\]
be the set of obstruction pairs. Suppose candidate field measurements $z_j$ have costs $c_j$.

\begin{proposition}[Monitoring realization]
A selected measurement set $S$ is sufficient to recover the repaired classification if and only if every pair in $E$ is separated by at least one $z_j$ with $j\in S$. Minimum-cost realization is therefore the weighted set-cover/test-cover problem
\[
\min_{S}\sum_{j\in S}c_j
\quad\text{subject to every pair in }E\text{ being separated.}
\]
\end{proposition}

This does not claim a new set-cover algorithm. It translates the formal witness set into the measurement requirement that an ecological monitoring design must satisfy.

Finally, the management consequence depends on utilities and state frequencies rather than partition size alone.

\begin{proposition}[No universal defect--regret ordering]
Without assumptions on target-state probabilities and the scale of management consequences, structural defect $\Delta_{\#}$ has no nontrivial universal monotone relationship with decision regret.
\end{proposition}

Reward rescaling changes regret without changing any partition, while arbitrarily many repaired states can have arbitrarily small decision consequences when their utility differences are small. Structural defect should therefore be interpreted as representation complexity only. Decision regret and value of information are separate quantities.

\subsection{Route coherence and minimum historical context}
Consider a rooted directed acyclic graph of declared replacement relations. Every root-to-terminal route carries the source label map to the terminal system.

\begin{theorem}[Route coherence]
A single route-independent carried classification, repaired classification, and structural defect exist if and only if every root-to-terminal route induces the same terminal carried label map.
\end{theorem}

\begin{theorem}[Minimum carried-map context]
When terminal carried maps disagree, one immutable context per distinct carried terminal map is necessary and sufficient to preserve all declared routes before exact refinement.
\end{theorem}

These results concern the semantics carried by externally declared replacement routes. They are not claims that history-sensitive equivalence is new in general \citep{MontanariPistore1997}. Their ecological interpretation is narrower: disturbance or management history should be retained only when alternative declared routes assign incompatible operational meanings to the same terminal configurations.

\section{Worked ecological example: restoration priority after pollinator turnover}
Consider two degraded sites, A and B, containing the same focal plant and showing the same coarse pollination output. Under the source monitoring and action repertoire, both are classified as ``pollination maintained,'' while a third state represents failure of pollination service. In the verified finite witness, the inherited target labels are $(0,0,1)$.

Structural change alters the pollinator community. The original dominant pollinator is lost, but one target configuration retains access to a substitute guild whereas the other does not. Competitor removal is introduced as a target intervention that can open floral access around the focal plant. Under this new action, the two states carrying the same inherited label have different successors: Site A lacks the substitute response channel, whereas Site B retains it. The inherited classification is therefore not decision-sufficient for the new intervention.

The local obstruction is exactly the A--B pair. Applying standard refinement to the inherited labels splits only that exposed fiber, changing $(0,0,1)$ to $(0,1,2)$. The structural defect is one macrostate. The ecological content of the repair is not ``add one arbitrary state'': it is ``record substitute-pollinator response capacity,'' because that is the distinction that explains the different action successors. Figure~\ref{fig:local-split} summarizes the audit-to-monitoring workflow.

\begin{figure}[t]
\centering
\includegraphics[width=0.98\linewidth]{../artifacts/figure1_local_split.svg}
\caption{Decision-sufficiency audit for the plant--pollinator example. Pollinator turnover leaves an inherited class in place, a target intervention exposes different successors within that class, and established refinement returns the least state distinction that a monitoring design must recover}
\label{fig:local-split}
\end{figure}

To make the management consequence explicit without presenting invented values as data, we add an illustrative one-step decision layer. Let intervention success at Site $i$ have probability $p_i$, success have benefit $B$, and intervention cost be $c_i$. The repaired model prefers Site B exactly when
\[
Bp_B-c_B>Bp_A-c_A,
\qquad\text{or equivalently}\qquad
p_B-p_A>\frac{c_B-c_A}{B}.
\]
An inherited model that pools A and B assigns the same success probability to both and therefore chooses the cheaper site whenever $c_A<c_B$.

For illustration, set $B=1$, $c_A=0.20$, $c_B=0.35$, $p_A=0.25$, and $p_B=0.80$. The inherited pooled probability is $\bar p=0.525$, giving expected net benefits $0.325$ for A and $0.175$ for B, so the inherited rule chooses A. The repaired model gives $0.05$ for A and $0.45$ for B, so it chooses B. Evaluated under the repaired consequences, using the inherited choice incurs one-step regret $0.40$. These numbers are illustrative, not empirical estimates; their purpose is to show which quantities a field application would need to estimate and the parameter region in which a structural split changes the management decision.

\begin{table}[t]
\centering
\small
\begin{tabular}{lccccc}
\toprule
Site & pooled $p$ & repaired $p_i$ & cost $c_i$ & pooled net benefit & repaired net benefit \\
\midrule
A & 0.525 & 0.25 & 0.20 & 0.325 & 0.05 \\
B & 0.525 & 0.80 & 0.35 & 0.175 & 0.45 \\
\bottomrule
\end{tabular}
\caption{Illustrative decision layer for the plant--pollinator example. The pooled inherited state favors cheaper Site A, whereas the repaired state information favors Site B}
\label{tab:decision}
\end{table}

\begin{figure}[t]
\centering
\includegraphics[width=0.78\linewidth]{../artifacts/figure4_decision_reversal.svg}
\caption{Region in which the repaired two-site decision prefers Site B. The boundary is $p_B-p_A=(c_B-c_A)/B$; the marked point is the illustrative case in Table~\ref{tab:decision}, not an empirical estimate}
\label{fig:decision-reversal}
\end{figure}

\begin{figure}[t]
\centering
\includegraphics[width=0.9\linewidth]{../artifacts/figure2_transport_defect.svg}
\caption{Structural repair complexity in the verified sharp finite family with independently exposed target distinctions. The curve concerns state-space compression only and is not a claim about monitoring cost or decision regret}
\label{fig:defect}
\end{figure}

The same example also illustrates route dependence. Suppose a terminal community can be reached either by pollinator loss followed by substitute colonization or by habitat modification followed by loss of the original guild. If both declared routes carry the same terminal labels, route identity is unnecessary for this interface. If they carry incompatible labels, no route-free inherited state assignment exists and an immutable carried-map context must be retained.

\begin{figure}[t]
\centering
\includegraphics[width=0.95\linewidth]{../artifacts/figure3_history_completion.svg}
\caption{Route coherence for carried ecological state meanings. Coherent routes share one carried interface; incompatible terminal carried maps require distinct immutable contexts before exact refinement}
\label{fig:history}
\end{figure}

\section{Discussion}
The main result is a change in where model transfer is audited. Ecological workflows commonly ask whether a model can be refitted, extrapolated, or updated after environmental change. Here the prior question is whether the state variable supplied to that workflow remains sufficient for the target decision. Structural change or a new intervention can expose a distinction that an inherited class was explicitly designed to ignore. When that happens, better parameter estimates within the old state representation do not solve the problem.

This perspective connects formal state abstraction to ecological monitoring without claiming that either literature begins here. Empirical state-and-transition work already shows how a small set of field variables can be selected to distinguish states and detect transitions \citep{JonesEtAl2023}. Our audit supplies a complementary trigger for redesign: if two configurations within one inherited state have different target action consequences, they must no longer be observationally interchangeable for that management problem. The obstruction pairs define what a candidate indicator set must separate. Ecological data are then needed to identify feasible indicators, estimate thresholds, and assess classification error.

The distinction from value of information is similarly important. VoI asks whether learning about an uncertain state, parameter, or model will improve expected management outcomes enough to justify sampling \citep{CanessaEtAl2015}. Our audit asks whether the current state representation is structurally capable of expressing the action consequence in the first place. The two approaches are sequential rather than competing: structural audit can identify a missing distinction; monitoring design can propose measurements that resolve it; and VoI can then evaluate whether collecting those measurements is worth the cost under uncertainty.

Separating structural defect from decision regret prevents an attractive but unsupported interpretation. A one-state split can be management-critical if it reverses a high-stakes intervention, while a large state-count expansion can be nearly irrelevant if the exposed distinctions are rare or have negligible consequences. $\Delta_{\#}$ and $\Delta_K$ therefore describe loss of compression, not ecological importance. Distribution-sensitive repair information, field measurement cost, and decision regret require additional assumptions and should be reported separately.

Historical context should be treated with the same economy. Ecology contains many legacy effects, but a model need not record a complete disturbance history simply because history can matter. The route-coherence criterion asks a more specific question: do alternative declared histories change the inherited operational meaning assigned to the present configuration? If not, history can be discarded for this decision interface. If so, the minimum completion retains only the distinct carried-map contexts required to make the present state unambiguous.

\subsection{Scope and limitations}
The framework is appropriate when investigators can declare a finite or finitely represented source system, a target system, the outputs and interventions that matter, and a defensible relation carrying source configurations to target configurations. Concrete uses could include auditing restoration-priority classes after interaction turnover, testing whether habitat or resilience categories remain decision-sufficient after disturbance, or checking whether an inherited community-state interface can support a newly available intervention.

It should not be used as a substitute for estimating transition probabilities, detecting species, inferring causal effects, learning the source--target relation, optimizing a policy, or quantifying observation and model uncertainty. Exactness is a benchmark for isolating structural failure. Approximate or stochastic portability, imperfect detection, and uncertainty in the structural relation are extensions rather than results established here.

The plant--pollinator case is a finite theoretical example. Its deterministic split is verified by the repository witnesses, and its probability/cost layer is an illustrative sensitivity calculation. It is not evidence that a particular plant, pollinator guild, or restoration program exhibits the stated probabilities or priority reversal.

The central ecological implication is therefore simple: a management state is not permanently valid merely because its label remains interpretable. Its adequacy depends on the distinctions exposed by the current system, observation target, action repertoire, and management objective. When structural change makes an inherited state insufficient, the useful output is not another unconstrained classification but a diagnosis of the missing distinction and a principled target for monitoring repair.

\section{Methods and proof organization}
The portability criterion follows by necessity and sufficiency of output, legal-row, and successor uniformity within carried fibers. Local obstruction is the corresponding finite counter-witness. The source-relative repair uses established monotone partition refinement \citep{PaigeTarjan1987,GivanDeanGreig2003}. The conditional-information inequality follows because every exact refinement of the carried classification refines $Q^\ast$. Monitoring realization is the covering problem over obstruction pairs. Non-monotonicity between structural defect and regret follows by changing utility scale without changing the partition. Route coherence is equality of complete terminal carried maps across declared routes; minimum context follows because distinct carried maps cannot share one immutable route-free label assignment. Detailed proof dependencies, finite witnesses, and reproducibility commands are provided in the supplementary manuscript and repository workflows.

\bibliographystyle{plainnat}
\bibliography{references}

\end{document}
